\documentclass[10pt,a4paper]{amsart}
\usepackage[margin=19mm]{geometry}
\usepackage{amsmath,amssymb,mathtools}
\usepackage[scr=rsfs,cal=euler]{mathalfa}
\usepackage{xfrac}
\usepackage[unicode,hidelinks,bookmarks=false]{hyperref}
\newcommand{\ie}{i.e.\ }
\newcommand{\cf}{cf.\ }
\newcommand{\resp}{respectively\ }
\newcommand{\Subsection}{\subsection}
\providecommand{\nobreakdash}{\nobreak\hbox{-}}
\setlength{\emergencystretch}{2em}
\multlinegap0pt
\usepackage{bm}
\usepackage{bbm}
\usepackage{dsfont}
\usepackage{xspace}
\usepackage{cancel}
\usepackage{mathtools}
\usepackage{amsthm}
\makeatletter
\@ifundefined{lem}{\newtheorem{lem}{Lemma}}{}
\@ifundefined{prop}{\newtheorem{prop}{Proposition}}{}
\makeatother
\newcommand{\Z}{\mathbb{Z}}
\title[New invariant surface measures for the cubic Schr\"odinger e]{New invariant surface measures for the cubic Schr\"odinger equation}
\author{Jean-Baptiste Casteras, Ana Bela Cruzeiro, Annie Millet}
\date{Source: arXiv:2502.02162v2 (CC BY 4.0)}
\begin{document}
\maketitle

\noindent\emph{Source and licence.} New invariant surface measures for the cubic Schr\"odinger equation. Version arXiv:2502.02162v2 (CC BY 4.0), distributed under the Creative Commons Attribution 4.0 International licence, as recorded in the arXiv licence metadata for this exact version and shown on the abstract page https://arxiv.org/abs/2502.02162v2. Full rights notice: licenses/arxiv-2502.02162-CC-BY-4.0.txt.\par\smallskip

\noindent\textit{Editorial scope.} Counted unit: the Prop whose source label is \texttt{None} in the article (source0.tex lines 706--709, source label \texttt{None}), delivered together with the article's own complete original proof of it (source0.tex lines 710--796, one \texttt{proof} environment of 87 lines). No mathematics is written, repaired or re-derived: statement and proof are the article's own text line for line. Required context retained uncounted: formvar (lines 420--424); :B: (lines 530--532); convseries (lines 1150--1160); convseries1 (lines 1153--1155). Every definition, display and label of those lines is retained unchanged. Omitted from this package and not redistributed: 1--419, 425--529, 533--705, 797--1149, 1156--1563. Nothing inside a retained excerpt is omitted or altered: every retained body is delivered exactly as the article's lines are written, so no omission receipt of any kind accompanies this package. Citations: the retained excerpts carry no citation command at all, so no bibliography entry of the article is transcribed and no reference is redistributed. Reference adaptation: none was needed and none was made; every reference inside the delivered unit resolves inside it, so no unresolved reference or citation is produced. Printed slips kept literal: no printed word, symbol, hypothesis or number of the counted statement or of its proof was altered. Layout and class: the article's own class is \texttt{amsart} with packages including inputenc, fontenc, bbm, amsmath, amsfonts, amssymb, amsthm, latexsym, xy, calc, multicol, color, appendix, srcltx; the wrapper is the lane's pinned \texttt{amsart} editorial wrapper at 10pt on A4, which restates the theorem-like environments the retained text needs and adds the article's own macro definitions that the retained text needs (\texttt{\textbackslash Z}), with extra wrapper packages (bm, bbm, dsfont, xspace, cancel, mathtools). The delivered numbering is the unit's own. Assets: no retained span contains a figure environment, a graphics-inclusion command, a PDF-page inclusion, a table or a TikZ picture; no image file is copied into this package, the package declares no asset dependency and no image file is redistributed with it. Licence: version arXiv:2502.02162v2 of this article is distributed under the Creative Commons Attribution 4.0 International licence, as recorded in the arXiv licence metadata for this exact version and shown on the abstract page https://arxiv.org/abs/2502.02162v2. A full rights notice is delivered at licenses/arxiv-2502.02162-CC-BY-4.0.txt. Characters: every retained excerpt is pure ASCII, so the delivered engine, XeLaTeX with its default Latin Modern text font, reports no missing character.
\begin{equation}
\label{formvar}
\mathbb{E}_{\mu_a} (|\varphi_k|^{2p})=\dfrac{c_p}{|k|^{ap}}, \quad \mbox{\rm where} \quad 
c_p= 2^p p!.
\end{equation} 

\begin{equation} 	\label{:B:}
(:B:)_k^N= B_k - 4 \; \Big(\sum_{m:1\leq |m| \leq |N|} \frac{1}{|m|^a}\Big) \varphi_k,
\end{equation}

\begin{lem}
\label{convseries} Let $d=1,2$. 
For any $\beta<0$, it holds that
\begin{equation} \label{convseries1}
\sum_{0\neq l,m\in \Z^2} \dfrac{1}{|l|^2}\dfrac{1}{|m|^2} \sum_{k\not\in \{0, l+m\}} \dfrac{|k|^{\beta}}{|l+m-k|^2}<\infty, 
\end{equation}
and
\begin{equation}  \label{convseries2}
\sum_{k\neq 0, m\not\in \{k, k/2\} } \dfrac{|k|^{\beta}}{|k-m|^4 |k-2m|^2}<\infty.
\end{equation}
\end{lem}

\begin{prop}
 Let $d=1,2$. For $a=2$ and $\beta <0$, the gradient of the vector field $B$ for $d=1$ (resp. $:B:$ for $d=2$) belongs to $L^{2p}_{\mu_a} (H^\beta ; H.S. (H^1 ,H^\beta))$ for any integer $p\geq 1$.

\end{prop}

\begin{proof}
As in the previous proof, we work at fixed $N$ but we drop the subscript $N$ in $(:B:)^N$.
A straight-forward computation gives  
\begin{align*}
D_{e_j} B_k (\varphi)&
= \sum_{l,m} (\varphi_{l+m} \overline{\varphi_l} \delta_{j,k-m}+ \varphi_{k-m} \overline{\varphi_l} \delta_{j,l+m}+ \varphi_{l+m} \varphi_{k-m} \delta_{-j,l}) \\
&=  \sum_{l\notin \{ 0, j-k\}} \varphi_{l+k-j} ( 2\overline{\varphi_l} + \varphi_{-l}).  
\end{align*}
Thus, we have
$$|D_{e_j} B_k (\varphi)|^2 = \sum_{l,l^\prime \notin \{ 0, j-k\}} \varphi_{l+k-j} (2 \overline{\varphi_l}+ \varphi_l )  \overline{\varphi_{l^\prime+k-j}} (2\varphi_{l^\prime}+ \overline{\varphi_{l^\prime}}) .$$

Let us first consider  the $1$-dimensional case. Since
$$\|\nabla B(\varphi) (\hat{e}_j)\|_{H^\beta}^2 \leq c \sum_{k} \dfrac{|k|^{2\beta}}{|j|^{2}}\;  \Big| \sum_{l\notin \{ 0, j-k\}}
 \varphi_{l+k-j}  (2 \overline{\varphi_l} + \varphi_{-l})  \Big|^2, $$
 we deduce
\begin{align*}
E_{\mu_a} \big[ \|\nabla B (\varphi)\|^2_{H.S. (H^1 ,H^\beta) }\big] &\leq  c \sum_{k,j} \dfrac{|k|^{2\beta}}{|j|^{2}} \;  E_{\mu_a} \Big[ \sum_{l,l'\notin \{0,  j-k\}}   T_{j,k}(l,l')\Big],
\end{align*}
 where $T_{j,k}(l,l')=\varphi_{l+k-j}  (2\overline{\varphi_l} +\varphi_{-l}) \overline{\varphi_{l'+k-j}} (2 \varphi_{l'} + \overline{\varphi_{-l'}}) $.
\bigskip

\noindent {\it Case 1:} $l=l'$, $T_{j,k}(l,l) = |\varphi_{l+k-j}|^2 (4 |\varphi_l|^2 + 2 \overline{\varphi_{l}} \overline{\varphi_{-l}} + 2 \varphi_{-l} \varphi_l  + |\varphi_{-l}|^2)$. 

\;  $\bullet$ If $k=j$, $E_{\mu_a}[T_{j,j}(l,l) ]=4  E_{\mu_a}[  |\varphi_l|^4 ] + E_{\mu_a} [|\varphi_l|^2 ] E_{\mu_a} [ |\varphi_{-l}|^2]$. 

\;  $\bullet$ If $k\neq j$ and $2l+k-j\neq 0$, $E_{\mu_a}[T_{j,k}(l,l) ] = 4 E_{\mu_a}[|\varphi_{l+k-j}|^2] E_{\mu_a}[ |\varphi_l|^2 ] + E_{\mu_a}[|\varphi_{l+k-j}|^2]  E_{\mu_a} [|\varphi_{-l}|^2] $. 

\;  $\bullet$ If  $k\neq j$ and $2l+k-j=  0$, $E_{\mu_a}[T_{j,k}(l,l) ]=  4 E_{\mu_a}[|\varphi_{l+k-j}|^2] E_{\mu_a}[ |\varphi_l|^2 ]  + E_{\mu_a}[|\varphi_{l+k-j}|^4]$.  
\bigskip

\noindent{\it Case 2:} $l\neq l'$. 

\indent $\bullet$ If $k=j$,  $T_{j,j}(l,l') = 4  |\varphi_l|^2   |\varphi_{l'}|^2 + 2 |\varphi_{l'}|^2 \varphi_l  \varphi_{-l}  + 2 |\varphi_l|^2 \overline{\varphi_{l'}} \overline {\varphi_{-l'}}
+ \overline{ \varphi_{l'}} \overline{ \varphi_{-l'}}$.  Therefore,
$E_{\mu_a}[ T_{j,j}(l,l') ] = 4  E_{\mu_a} [ |\varphi_l|^2]  E_{\mu_a}[  |\varphi_{l'}|^2]$.  

\indent $\bullet$ If $k\neq j$, $E_{\mu_a}[T_{j,k}(l,l')]=0$, except if $ l+k-j=-l'$. In that case we have
$E_{\mu_a}[T_{j,k}(l,l')] = E_{\mu_a} [|\varphi_{l+k-j}|^2] E_{\mu_a} [|\varphi_{-l}|^2]$ if $l+k-j\neq -l$, and $E_{\mu_a} [|\varphi_{l+k-j}|^4]$ if
$l+k-j=  -l$. 
\bigskip

 For $a=2$ and $\beta <0$, using once more \eqref{convseries1}, 
  we have 
%\[ \sum_{k,j\neq 0} \dfrac{|k|^{2\beta}}{|j|^{2}} \Big[ \sum_{l\neq 0} |l|^{-2a} + \sum_{l\notin \{0, j-k\}} |l|^{-a}|l+k-j|^{-a} +
%\sum_{l,l'\neq 0, l\neq l'} |l|^{-a} |l'|^{-a}\Big] <\infty.\]
\begin{align*}
\sum_{k\neq 0} &\frac{|k|^{2\beta}}{|k|^2} \Big[ \sum_{l\neq 0} |l|^{-2a} + \sum_{l,l' \neq 0, l\neq l'} |l|^{-a} |l'|^{-a} \Big] \\
& + 
 \sum_{k,j\neq 0, k\neq j} \dfrac{|k|^{2\beta}}{|j|^{2}}  \sum_{l\notin \{0, j-k\}}\big(  |l|^{-a}|l+k-j|^{-a} +
 |l+k-j|^{-2a}\big) <\infty.
\end{align*} 



Next, we deal with the $2$-dimensional case. We need to be more careful in this case and we begin by computing $\|D_{e_j} B_k (\varphi)\|_{L^2_{\mu_a}}$. As in the one-dimensional case, we have
\begin{align*}
\int & |D_{e_j}B_k(\varphi)|^2 d\mu_a =   \delta_{j,k} \Big(  4 \sum_{l\neq 0}  E_{\mu_a}[|\varphi_l|^4] + \sum_{l\neq 0} E_{\mu_a}[|\varphi_l|^2]  E_{\mu_a}[|\varphi_{-l}|^2]\\
& \qquad \qquad + 4 \sum_{l,l'\neq 0, l\neq l'} E_{\mu_a} [|\varphi_l|^2] E_{\mu_a} [|\varphi_{l'}|^2|] \Big)\\
& + 1_{\{j\neq k\} } \Big( 4 \sum_{l  \notin \{0, j-k\} } E_{\mu_a}[ |\varphi_{l+k-j}|^2]  E_{\mu_a}[ |\varphi_l|^2]  +2 1_{\{ j-k \; {\rm even}\}} E_{\mu_a} [ |\varphi_{(k-j)/2}|^4]  \\
& \qquad \qquad +2  { \sum_{l,2l \notin \{0, j-k\}} } E_{\mu_a}[ |\varphi_{l+k-j}|^2]  E_{\mu_a}[ |\varphi_{-l}|^2] \Big).
\end{align*}
Let 
\begin{align*}
 {\mathcal J}_N&(j,k)= \delta_{j,k} \Big(  4 \sum_{l\neq 0}  E_{\mu_a}[|\varphi_l|^4] + \sum_{l\neq 0} E_{\mu_a}[|\varphi_l|^2]  E_{\mu_a}[|\varphi_{-l}|^2] \Big) \\
&+  1_{\{j\neq k\} }\Big(  4 \sum_{l  \notin \{0, j-k\} } E_{\mu_a}[ |\varphi_{l+k-j}|^2]  E_{\mu_a}[ |\varphi_l|^2]  +2 1_{\{ j-k \; {\rm even}\}} E_{\mu_a} [ |\varphi_{(k-j)/2}|^4]  \\
& \qquad \qquad +2 {  \sum_{l,2l \notin \{0, j-k\}} } E_{\mu_a}[ |\varphi_{l+k-j}|^2]  E_{\mu_a}[ |\varphi_{-l}|^2] \Big).
\end{align*}


Using \eqref{:B:}, we deduce that 
$$D_{e_j} (:B :)_k (\varphi) = D_{e_j} B_k (\varphi) - C_N \delta_{j,k} .$$
Therefore, since $C_N=2 D_N$, 
\begin{align*}
\int & |D_{e_j}B_k(\varphi)|^2 d\mu_a =   {\mathcal J}_N(j,k) + \delta_{j,k}\big( 4 D_N^2 - 2C_N \, 2D_N + C_N^2\big)  = {\mathcal J}_N(j,k) . 
\end{align*}
Using \eqref{formvar} and \eqref{convseries1} in Lemma \ref{convseries}, we deduce that
$$\sup_N \sum_{j,k} \frac{|k|^{2\beta}}{|j|^2} {\mathcal J}_N(j,k) <\infty.$$
For more general integers $p$, we proceed recursively. Indeed, we have formally (as previously, we use erroneously the independence
 of $ |D_{e_j}B_k(\varphi)|^2$ and $|D_{e_j}B_k(\varphi)|^{2p}$, but the argument can be made rigorous)  
\begin{align*}
E_{\mu_a}\big(  |D_{e_j}B_k(\varphi)|^{2p+2}\big) 
&=  E_{\mu_a} \big( |D_{e_j}B_k(\varphi)|^{2p} |D_{e_j}B_k(\varphi)|^2 \big)\\
&=  E_{\mu_a} \big( |D_{e_j}B_k(\varphi)|^{2p}  {\mathcal J}_N(j,k)  \big)\\
&\leq C   E_{\mu_a} \big( |D_{e_j}B_k(\varphi) |^{2p} \big).
\end{align*}
This concludes the proof.
\end{proof}

\end{document}
